\documentclass[aspectratio=169]{beamer}

% Thème simple et propre
\usetheme{Madrid}
\usecolortheme{default}

% Encodage
\usepackage[T1]{fontenc}
\usepackage[utf8]{inputenc}
\usepackage[french]{babel}

% Maths et images
\usepackage{amsmath, amssymb}
\usepackage{graphicx}
\usepackage{booktabs}

% Enlever les icônes de navigation
\setbeamertemplate{navigation symbols}{}

% Infos du document
\title{Coursidor \& Robozzle}
\author{Jontef Raphaël}
\institute{Enseirb-Matmeca}
\date{Semestre 6, année 2025-2026}

\begin{document}

% -------------------------------------------------
\begin{frame}{Pureté et impureté}

\begin{columns}[T]

% -------- COLONNE GAUCHE --------
\begin{column}{0.40\textwidth}

\textbf{Coursidor est impur, Robozzle est pur}
\begin{itemize}
    \item Coursidor : état global, effets de bord
    \item Robozzle : fonctions pures, objets souverains (\texttt{PuzzleState}) nouvel état à chaque opération
\end{itemize}

\vspace{1cm}
\textbf{Avantages de la pureté}
\begin{itemize}
    \item Séparation des comportements au sein des fonctions
    \item Facilité de test et de débogage
\end{itemize}
\end{column}

\begin{column}{0.40\textwidth}

\textbf{Nuances}
\begin{itemize}
    \item Coursidor : copie de petites structures comme \texttt{move\_t} ou \texttt{dir\_t}
    \item Robozzle : effets de bord sur le \texttt{DOM} en \texttt{HTML}, impureté intrinsèque à l'interface utilisateur
\end{itemize}

\vspace{0.4cm}

\textbf{Désavantages de la pureté}
\begin{itemize}
    \item Paradigme de niveau élevé, plus difficile à maîtriser
    \item Coût en mémoire très élevé
\end{itemize}

\end{column}




\end{columns}

\end{frame}

% -------------------------------------------------
\end{document}